\title{Επικουρικές Υπηρεσίες μη Σχετικές με τη Συχνότητα}

\section{Εισαγωγή}
\label{kef-02-theoria}

Η συχνότητα είναι κοινή σε ολόκληρη τη συγχρονισμένη περιοχή και εκφράζει το ισοζύγιο
ενεργής ισχύος του συστήματος. Οι υπηρεσίες που τη διατηρούν εξετάζονται στο
\hyperref[kef-03-theoria]{Κεφάλαιο 3}. Η ασφαλής λειτουργία απαιτεί επιπλέον υπηρεσίες που
αφορούν κυρίως την τάση. Σε αντίθεση με τη συχνότητα, το μέτρο της τάσης διαφέρει από ζυγό σε
ζυγό και καθορίζεται κυρίως από το ισοζύγιο άεργης ισχύος στην περιοχή του ζυγού. Οι
υπηρεσίες αυτές ονομάζονται \emph{επικουρικές υπηρεσίες μη σχετικές με τη συχνότητα}
(non-frequency ancillary services) και έχουν κατά κανόνα τοπικό χαρακτήρα: παρέχονται από
μέσα που βρίσκονται ηλεκτρικά κοντά στο σημείο όπου εμφανίζεται η ανάγκη.

Η Οδηγία (ΕΕ) 2019/944 για την εσωτερική αγορά ηλεκτρικής ενέργειας ορίζει ως επικουρική
υπηρεσία κάθε υπηρεσία απαραίτητη για τη λειτουργία συστήματος μεταφοράς ή διανομής. Οι
επικουρικές υπηρεσίες περιλαμβάνουν την εξισορρόπηση και τις υπηρεσίες μη σχετικές με τη
συχνότητα, όχι όμως τη διαχείριση συμφόρησης. Υπηρεσίες μη σχετικές με τη συχνότητα είναι
όσες χρησιμοποιεί ο Διαχειριστής Συστήματος Μεταφοράς (ΔΣΜ) ή ο Διαχειριστής Συστήματος
Διανομής (ΔΣΔ) για:

\begin{itemize}
\item τη ρύθμιση τάσης σε μόνιμη κατάσταση (steady state voltage control),
\item τη γρήγορη έγχυση άεργου ρεύματος (fast reactive current injection),
\item την αδράνεια για την τοπική ευστάθεια του δικτύου (inertia for local grid stability),
\item το ρεύμα βραχυκύκλωσης (short-circuit current),
\item την ικανότητα επανεκκίνησης μετά από γενική διακοπή (black start capability) και
\item την ικανότητα λειτουργίας σε νησίδα (island operation capability).
\end{itemize}

Στο παραδοσιακό σύστημα οι υπηρεσίες αυτές παρέχονταν σε μεγάλο βαθμό εγγενώς από τις
σύγχρονες γεννήτριες των μεγάλων σταθμών. Ο αυτόματος ρυθμιστής τάσης κάθε γεννήτριας
ρυθμίζει την τάση στους ακροδέκτες της, η σύγχρονη μηχανή τροφοδοτεί κατά τα βραχυκυκλώματα
ρεύμα πολλαπλάσιο του ονομαστικού και οι στρεφόμενες μάζες της αποθηκεύουν κινητική ενέργεια.
Η αντικατάσταση των μονάδων αυτών από αιολικά και φωτοβολταϊκά πάρκα, τα οποία συνδέονται
μέσω ηλεκτρονικών μετατροπέων ισχύος, περιορίζει τις εγγενείς αυτές συνεισφορές. Οι
μετατροπείς μπορούν να παρέχουν αντίστοιχες υπηρεσίες, εφόσον σχεδιαστούν και ρυθμιστούν
κατάλληλα, αλλά με διαφορετικά χαρακτηριστικά: το ρεύμα τους δεν υπερβαίνει πολύ το ονομαστικό
και η απόκρισή τους καθορίζεται εξ ολοκλήρου από τον έλεγχό τους. Για τον λόγο αυτό ο
Κανονισμός (ΕΕ) 2016/631 για τις απαιτήσεις σύνδεσης των μονάδων ηλεκτροπαραγωγής
(Requirements for Generators, RfG) ορίζει ικανότητες άεργης ισχύος, ρύθμισης τάσης και
γρήγορης έγχυσης ρεύματος σφάλματος, ενώ οι ΔΣΜ εγκαθιστούν επιπλέον εξοπλισμό, όπως πηνία
και πυκνωτές αντιστάθμισης, διατάξεις FACTS και σύγχρονους αντισταθμιστές. Ο
Πίνακας~\ref{tab:nfas} συνοψίζει τον σκοπό και τη χρονική κλίμακα κάθε υπηρεσίας.

\begin{table}[!htbp]
\centering
\begin{tabular}{lll}
\hline
Υπηρεσία & Σκοπός & Χρονική κλίμακα \\
\hline
ρύθμιση τάσης & τάσεις εντός ορίων & s -- h \\
έγχυση άεργου ρεύματος & στήριξη της τάσης στα σφάλματα & ms -- s \\
αδράνεια & αρχική απόκριση της συχνότητας & ms -- s \\
ρεύμα βραχυκύκλωσης & λειτουργία της προστασίας & ms \\
επανεκκίνηση & αποκατάσταση μετά από διακοπή & min -- h \\
λειτουργία σε νησίδα & τροφοδότηση απομονωμένου τμήματος & min -- h \\
\hline
\end{tabular}
\caption{Επικουρικές υπηρεσίες μη σχετικές με τη συχνότητα: σκοπός και χρονική κλίμακα.}
\label{tab:nfas}
\end{table}

Το κεφάλαιο δίνει έμφαση στις δύο πρώτες υπηρεσίες, οι οποίες αποτελούν τη βάση της στήριξης
της τάσης. Αρχικά εξετάζονται η σχέση της τάσης με την άεργη ισχύ, οι γραμμές μεταφοράς ως
πηγές και καταναλωτές άεργης ισχύος, τα μέσα ρύθμισης της τάσης, η ευστάθεια τάσης και η
ιεραρχική οργάνωση της ρύθμισης. Ακολουθούν οι διατάξεις FACTS και η γρήγορη έγχυση άεργου
ρεύματος κατά τα σφάλματα, ενώ το ρεύμα βραχυκύκλωσης, η επανεκκίνηση και η λειτουργία σε
νησίδα εξετάζονται συνοπτικά. Η αδράνεια εξετάζεται μαζί με τις υπηρεσίες συχνότητας στο
\hyperref[kef-03-theoria]{Κεφάλαιο 3}.

\section{Ρύθμιση τάσης σε μόνιμη κατάσταση}
\label{sec:v-control}

Σκοπός της ρύθμισης τάσης είναι να διατηρούνται οι τάσεις όλων των ζυγών εντός των
επιτρεπόμενων ορίων, τόσο σε κανονική λειτουργία όσο και μετά από απρόβλεπτα συμβάντα. Οι
υψηλές τάσεις καταπονούν τη μόνωση του εξοπλισμού. Οι χαμηλές τάσεις αυξάνουν τα ρεύματα και
τις απώλειες για την ίδια μεταφερόμενη ισχύ, μειώνουν τα περιθώρια ευστάθειας και προκαλούν
δυσλειτουργίες στον εξοπλισμό των καταναλωτών. Για την Ηπειρωτική Ευρώπη ο SO GL ορίζει ως
εύρος κανονικής λειτουργίας 0,90--1,05 p.u. για το επίπεδο των 400 kV, δηλαδή 360--420 kV,
και 0,90--1,118 p.u. για τάσεις από 110 έως 300 kV, όπως τα 150 kV του ΕΣΜΗΕ. Παράλληλα, η
ρύθμιση επιδιώκει τον περιορισμό των ροών άεργης ισχύος, ώστε να μειώνονται οι απώλειες και
να αποδεσμεύεται μεταφορική ικανότητα για την ενεργή ισχύ, καθώς και τη διατήρηση επαρκούς
εφεδρείας άεργης ισχύος για την αντιμετώπιση διαταραχών.

\subsection{Τάση και άεργη ισχύς}
\label{sec:v-q}

Η σχέση της τάσης με την άεργη ισχύ προκύπτει από τη γραμμή δύο ζυγών του
Σχήματος~\ref{fig:vq}, με σύνθετη αντίσταση σειράς $R + jX$, μέσω της οποίας ρέει ισχύς
$P + jQ$ προς τον ζυγό 2. Αν ο φασιθέτης $\tilde{V}_2 = V_2$ ληφθεί ως αναφορά, το ρεύμα της
γραμμής είναι $\tilde{I} = (P - jQ)/V_2$ και η τάση του ζυγού 1 είναι

\begin{equation}
\tilde{V}_1 = V_2 + (R + jX)\,\tilde{I}
= V_2 + \frac{RP + XQ}{V_2} + j\,\frac{XP - RQ}{V_2}
\label{eq:v1phasor}
\end{equation}

Ο δεύτερος όρος βρίσκεται σε φάση με τον $\tilde{V}_2$ και καθορίζει κυρίως τη διαφορά των
μέτρων των τάσεων, ενώ ο τρίτος είναι κάθετος και καθορίζει τη διαφορά των γωνιών
$\delta = \theta_1 - \theta_2$. Όταν η γωνία $\delta$ είναι μικρή, ισχύει κατά προσέγγιση

\begin{equation}
\Delta V = V_1 - V_2 \approx \frac{RP + XQ}{V_2}, \qquad
\delta \approx \frac{XP - RQ}{V_2^2}
\label{eq:dv}
\end{equation}

Στα δίκτυα μεταφοράς η αντίδραση των γραμμών είναι πολλαπλάσια της αντίστασής τους, οπότε
$\Delta V \approx XQ/V_2$ και $\delta \approx XP/V_2^2$: η διαφορά των μέτρων των τάσεων
καθορίζεται από τη ροή άεργης ισχύος και η διαφορά γωνιών από τη ροή ενεργής ισχύος. Για
$R = 0$ η \eqref{eq:v1phasor}, αναλυμένη σε συνιστώσες, δίνει τις ακριβείς σχέσεις

\begin{equation}
P = \frac{V_1 V_2}{X}\,\sin\delta, \qquad
Q = \frac{V_1 V_2 \cos\delta - V_2^2}{X}
\label{eq:pq-line}
\end{equation}

όπου $Q$ η άεργη ισχύς που φθάνει στον ζυγό 2. Στα δίκτυα διανομής ο λόγος $X/R$ είναι
μικρότερος και η ενεργή ισχύς επηρεάζει σημαντικά και το μέτρο της τάσης· το θέμα εξετάζεται
στο Κεφάλαιο 5.

\begin{figure}[!htbp]
\centering
\includegraphics[width=\linewidth]{figures/v-q-phasor.svg}
\caption{Γραμμή δύο ζυγών με αντίδραση σειράς $X$ (αριστερά) και διανυσματικό διάγραμμα των
τάσεων για ρεύμα που καθυστερεί της τάσης $V_2$ (δεξιά). Η πτώση τάσης $jX\tilde{I}$
αναλύεται σε μια συνιστώσα $XQ/V_2$ σε φάση με την $V_2$, η οποία καθορίζει κυρίως τη
διαφορά των μέτρων, και σε μια κάθετη συνιστώσα $XP/V_2$, η οποία καθορίζει τη διαφορά
γωνιών $\delta$.}
\label{fig:vq}
\end{figure}

Από τις σχέσεις αυτές προκύπτουν τρία βασικά χαρακτηριστικά της ρύθμισης τάσης:

\begin{itemize}
\item Η άεργη ισχύς ρέει από τον ζυγό με τη μεγαλύτερη προς τον ζυγό με τη μικρότερη τάση,
      όπως η ενεργή ισχύς ρέει από τον ζυγό με τη μεγαλύτερη προς τον ζυγό με τη μικρότερη
      γωνία.
\item Η μεταφορά άεργης ισχύος σε μεγάλες αποστάσεις δεν είναι αποδοτική. Απαιτεί μεγάλες
      διαφορές τάσης και αυξάνει το ρεύμα, άρα τις απώλειες ενεργής ισχύος $RI^2$ και την
      άεργη ισχύ $XI^2$ που καταναλώνει η ίδια η γραμμή. Η άεργη ισχύς πρέπει επομένως να
      παράγεται κοντά στο σημείο όπου καταναλώνεται.
\item Η τάση είναι τοπικό μέγεθος. Ενώ η συχνότητα ρυθμίζεται από οποιαδήποτε μονάδα της
      συγχρονισμένης περιοχής, η τάση ενός ζυγού ρυθμίζεται αποτελεσματικά μόνο από μέσα
      που βρίσκονται ηλεκτρικά κοντά του.
\end{itemize}

Η επίδραση μιας μεταβολής $\Delta Q$ της άεργης ισχύος που εγχέεται σε έναν ζυγό εκτιμάται με
το ισοδύναμο Thévenin του δικτύου, όπως φαίνεται από τον ζυγό: μια πηγή τάσης πίσω από
αντίδραση $X_{\mathrm{th}}$. Από την \eqref{eq:dv} για $R = 0$ προκύπτει
$\Delta V \approx X_{\mathrm{th}}\,\Delta Q / V$ και, με την ισχύ βραχυκύκλωσης
$S_{\mathrm{sc}} = V^2 / X_{\mathrm{th}}$ του ζυγού,

\begin{equation}
\frac{\Delta V}{V} \approx \frac{\Delta Q}{S_{\mathrm{sc}}}
\label{eq:dv-ssc}
\end{equation}

Η ευαισθησία της τάσης στην άεργη ισχύ είναι επομένως αντιστρόφως ανάλογη της ισχύος
βραχυκύκλωσης. Σε ζυγό με μεγάλη ισχύ βραχυκύκλωσης, δηλαδή ηλεκτρικά κοντά σε πολλές
σύγχρονες μηχανές, η τάση μεταβάλλεται ελάχιστα, ενώ σε ασθενή ζυγό μικρές μεταβολές της
άεργης ισχύος μεταβάλλουν σημαντικά την τάση. Η σχέση συνδέει τη ρύθμιση τάσης με το ρεύμα
βραχυκύκλωσης, το οποίο αποτελεί χωριστή υπηρεσία.

\emph{Παράδειγμα.} Συστοιχία πυκνωτών 50 Mvar συνδέεται σε ζυγό 150 kV με ισχύ
βραχυκύκλωσης 2\,500 MVA. Σύμφωνα με την \eqref{eq:dv-ssc}, η τάση αυξάνεται κατά
50/2\,500 = 2\%, δηλαδή κατά 3 kV περίπου. Η ίδια συστοιχία σε ζυγό 400 kV με ισχύ
βραχυκύκλωσης 10\,000 MVA αυξάνει την τάση μόλις κατά 0,5\%.

\subsection{Γραμμές μεταφοράς και φυσική ισχύς}
\label{sec:sil}

Στις γραμμές μεγάλου μήκους και στα καλώδια η εγκάρσια χωρητικότητα δεν μπορεί να αγνοηθεί.
Μια γραμμή \emph{παράγει} άεργη ισχύ στη χωρητικότητά της, ανάλογη του τετραγώνου της τάσης,
και \emph{καταναλώνει} άεργη ισχύ στην αυτεπαγωγή της, ανάλογη του τετραγώνου του ρεύματος.
Για γραμμή μήκους $\ell$, με αυτεπαγωγή $L$ και χωρητικότητα $C$ ανά μονάδα μήκους, και για
τάση και ρεύμα περίπου σταθερά κατά μήκος της, η καθαρή άεργη ισχύς που παράγει η γραμμή είναι

\begin{equation}
Q_{\mathrm{line}} \approx \omega \ell \left( C V^2 - L I^2 \right)
\label{eq:qline}
\end{equation}

Οι δύο όροι εξισώνονται όταν $V/I = \sqrt{L/C} = Z_c$, δηλαδή όταν η γραμμή φορτίζεται με την
\emph{κυματική} (χαρακτηριστική) σύνθετη αντίστασή της $Z_c$. Η αντίστοιχη ισχύς

\begin{equation}
P_0 = \frac{V_n^2}{Z_c}
\label{eq:sil}
\end{equation}

όπου $V_n$ η ονομαστική πολική τάση, ονομάζεται \emph{φυσική ισχύς} της γραμμής (Surge
Impedance Loading, SIL). Ανάλογα με τη μεταφερόμενη ισχύ διακρίνονται τρεις περιπτώσεις:

\begin{itemize}
\item Για $P = P_0$ η γραμμή ούτε παράγει ούτε καταναλώνει καθαρή άεργη ισχύ και η τάση είναι
      σταθερή κατά μήκος της.
\item Για $P < P_0$, όπως στις ώρες χαμηλού φορτίου, υπερισχύει η χωρητικότητα: η γραμμή
      παράγει άεργη ισχύ και η τάση αυξάνεται προς το μέσον της ή προς το ανοιχτό άκρο της.
\item Για $P > P_0$ υπερισχύει η αυτεπαγωγή: η γραμμή καταναλώνει άεργη ισχύ, η οποία πρέπει
      να καλυφθεί από τα άκρα της, και η τάση μειώνεται κατά μήκος της.
\end{itemize}

Για εναέριες γραμμές 400 kV με δέσμη αγωγών η $Z_c$ είναι περίπου 250--300 $\Omega$ και η
φυσική ισχύς περίπου 550--650 MW, ενώ για γραμμές 150 kV η $Z_c$ είναι περίπου 350--400
$\Omega$ και η φυσική ισχύς 55--65 MW. Τα καλώδια έχουν πολύ μεγαλύτερη χωρητικότητα και πολύ
μικρότερη $Z_c$, οπότε η φυσική τους ισχύς υπερβαίνει κατά πολύ τη θερμική τους ικανότητα.
Ένα καλώδιο εναλλασσόμενου ρεύματος παράγει επομένως πάντα καθαρή άεργη ισχύ, και οι μεγάλες
καλωδιακές συνδέσεις, όπως οι υποβρύχιες διασυνδέσεις νησιών, απαιτούν πηνία αντιστάθμισης.

Η κατανομή της τάσης κατά μήκος μιας γραμμής χωρίς απώλειες προκύπτει από τις εξισώσεις της
γραμμής μεγάλου μήκους. Η τάση σε απόσταση $x$ από το άκρο λήψης είναι

\begin{equation}
\tilde{V}(x) = \tilde{V}_r \cos\beta x + j Z_c \tilde{I}_r \sin\beta x
\label{eq:longline}
\end{equation}

όπου $\tilde{V}_r$ και $\tilde{I}_r$ η τάση και το ρεύμα στο άκρο λήψης και
$\beta = \omega\sqrt{LC}$ η σταθερά φάσης, περίπου 0,06° ανά km για εναέριες γραμμές στα
50 Hz. Όταν η γραμμή μεταφέρει τη φυσική της ισχύ, $\tilde{I}_r = \tilde{V}_r / Z_c$ και η
\eqref{eq:longline} δίνει $\tilde{V}(x) = \tilde{V}_r\, e^{j\beta x}$: το μέτρο της τάσης
είναι σταθερό και μεταβάλλεται μόνο η γωνία της. Όταν η γραμμή είναι ανοιχτή στο άκρο λήψης,
$\tilde{I}_r = 0$ και η τάση του ανοιχτού άκρου είναι

\begin{equation}
V_r = \frac{V_s}{\cos\beta\ell}
\label{eq:ferranti}
\end{equation}

όπου $V_s$ η τάση του άκρου τροφοδότησης. Επειδή $\cos\beta\ell < 1$, η τάση του ανοιχτού
άκρου υπερβαίνει την τάση τροφοδότησης. Η αύξηση αυτή ονομάζεται \emph{φαινόμενο Ferranti}.

\emph{Παράδειγμα.} Για εναέρια γραμμή μήκους 400 km είναι $\beta\ell \approx 24^\circ$ και,
σύμφωνα με την \eqref{eq:ferranti}, η τάση του ανοιχτού άκρου είναι περίπου 1,095 φορές η
τάση τροφοδότησης. Αν η γραμμή ηλεκτρίζεται από ζυγό 400 kV, η τάση στο ανοιχτό άκρο φθάνει
τα 438 kV, πάνω από το όριο των 420 kV. Για τον λόγο αυτό οι μεγάλες γραμμές υπερυψηλής
τάσης εξοπλίζονται με πηνία αντιστάθμισης στα άκρα τους.

Το Σχήμα~\ref{fig:sil} δείχνει την κατανομή της τάσης κατά μήκος της ίδιας γραμμής, με σταθερή
τάση 1 p.u. στα δύο άκρα, για διάφορες τιμές της μεταφερόμενης ισχύος, καθώς και για
τροφοδότηση από το ένα άκρο με το άλλο ανοιχτό.

\begin{figure}[!htbp]
\centering
\includegraphics[width=\linewidth]{figures/line-voltage-profile.svg}
\caption{Κατανομή της τάσης κατά μήκος εναέριας γραμμής 400 km χωρίς απώλειες. Οι συνεχείς
καμπύλες αντιστοιχούν σε τάση 1 p.u. και στα δύο άκρα: χωρίς φορτίο η τάση αυξάνεται προς το
μέσον, με φόρτιση ίση με τη φυσική ισχύ $P_0$ παραμένει σταθερή, ενώ με φόρτιση $2P_0$
μειώνεται. Η διακεκομμένη καμπύλη αντιστοιχεί σε γραμμή ανοιχτή στο δεξί άκρο, όπου η τάση
υπερβαίνει το όριο των 420 kV (1,05 p.u.) του επιπέδου των 400 kV.}
\label{fig:sil}
\end{figure}

Η φυσική ισχύς εξηγεί γιατί οι ανάγκες ρύθμισης τάσης μεταβάλλονται κατά τη διάρκεια της
ημέρας. Στις ώρες αιχμής οι γραμμές φορτίζονται πάνω από τη φυσική τους ισχύ και καταναλώνουν
άεργη ισχύ, οπότε οι τάσεις τείνουν να μειωθούν και απαιτείται παραγωγή άεργης ισχύος. Στις
ώρες χαμηλού φορτίου οι γραμμές φορτίζονται κάτω από τη φυσική τους ισχύ και παράγουν άεργη
ισχύ, οπότε οι τάσεις τείνουν να αυξηθούν και απαιτείται απορρόφηση άεργης ισχύος. Οι υψηλές
τάσεις είναι συχνότερες όταν η διεσπαρμένη παραγωγή στα δίκτυα διανομής καλύπτει μεγάλο μέρος
της ζήτησης, οπότε μειώνεται η ροή ισχύος στο δίκτυο μεταφοράς, καθώς και σε δίκτυα με πολλά
καλώδια.

\subsection{Μέσα ρύθμισης της τάσης}
\label{sec:q-sources}

Τα μέσα ρύθμισης της τάσης είτε παράγουν ή απορροφούν άεργη ισχύ είτε ανακατανέμουν τις ροές
της στο δίκτυο.

\textbf{Σύγχρονες γεννήτριες.} Ο αυτόματος ρυθμιστής τάσης (Automatic Voltage Regulator, AVR)
μεταβάλλει το ρεύμα διέγερσης της γεννήτριας, ώστε η τάση στους ακροδέκτες της, ή σε σημείο
του δικτύου κοντά σε αυτήν, να παραμένει στην τιμή αναφοράς. Όταν η γεννήτρια είναι
υπερδιεγερμένη παράγει άεργη ισχύ και όταν είναι υποδιεγερμένη απορροφά. Τα όρια της
λειτουργίας της δίνονται από το διάγραμμα ικανότητας (capability chart) του
Σχήματος~\ref{fig:capability}:

\begin{itemize}
\item το όριο ρεύματος στάτη, κύκλος σταθερής φαινόμενης ισχύος, λόγω της θέρμανσης του
      τυλίγματος του στάτη,
\item το όριο ρεύματος διέγερσης, λόγω της θέρμανσης του τυλίγματος του δρομέα, το οποίο
      περιορίζει την παραγωγή άεργης ισχύος,
\item το όριο υποδιέγερσης, λόγω ευστάθειας και θέρμανσης των άκρων του πυρήνα του στάτη, το
      οποίο περιορίζει την απορρόφηση άεργης ισχύος, και
\item το όριο της κινητήριας μηχανής, για την ενεργή ισχύ.
\end{itemize}

Όταν το ρεύμα διέγερσης υπερβεί το όριό του, ο περιοριστής υπερδιέγερσης (Over-Excitation
Limiter, OEL) το επαναφέρει μετά από χρονική καθυστέρηση στην επιτρεπόμενη τιμή. Η γεννήτρια
παύει τότε να ρυθμίζει την τάση και συμπεριφέρεται ως πηγή σταθερής άεργης ισχύος, γεγονός που,
όπως φαίνεται στην Ενότητα «\nameref{sec:v-stability}», συμβάλλει στην αστάθεια τάσης.

\textbf{Μονάδες που συνδέονται μέσω μετατροπέων.} Τα αιολικά και φωτοβολταϊκά πάρκα, οι
σταθμοί αποθήκευσης με συσσωρευτές και οι σταθμοί HVDC με μετατροπείς πηγής τάσης ρυθμίζουν
την άεργη ισχύ τους ανεξάρτητα από την ενεργή, εντός του ορίου ρεύματος του μετατροπέα,
$P^2 + Q^2 \le (V I_{\max})^2$. Σε αντίθεση με τη σύγχρονη γεννήτρια, η περιοχή λειτουργίας
είναι συμμετρική ως προς την παραγωγή και την απορρόφηση άεργης ισχύος και συρρικνώνεται
ανάλογα με την τάση (Σχήμα~\ref{fig:capability}). Ένας μετατροπέας με ονομαστική φαινόμενη
ισχύ ίση με τη μέγιστη ενεργή ισχύ της μονάδας δεν διαθέτει περιθώριο άεργης ισχύος όταν η
μονάδα αποδίδει τη μέγιστη ενεργή ισχύ της· γι' αυτό οι μετατροπείς υπερδιαστασιολογούνται.
Αντίθετα, όταν η ενεργή ισχύς είναι μηδενική, για παράδειγμα σε φωτοβολταϊκό σταθμό τη νύχτα,
ο μετατροπέας μπορεί να διαθέσει όλη την ονομαστική του ισχύ ως άεργη, εφόσον παραμένει
συνδεδεμένος, λειτουργώντας ουσιαστικά ως STATCOM. Η άεργη ισχύς ρυθμίζεται σε σταθερή τιμή,
με σταθερό συντελεστή ισχύος ή ως συνάρτηση της τάσης με στατισμό· μόνο ο τελευταίος τρόπος
αποκρίνεται αυτόματα στις μεταβολές της τάσης. Ο RfG ορίζει το πλαίσιο της απαιτούμενης
ικανότητας άεργης ισχύος και των τρόπων ρύθμισης, ενώ τις συγκεκριμένες τιμές καθορίζει ο
αρμόδιος διαχειριστής. Η συμμετοχή των μονάδων των δικτύων διανομής στη ρύθμιση της τάσης
εξετάζεται στο Κεφάλαιο 5.

\begin{figure}[!htbp]
\centering
\includegraphics[width=\linewidth]{figures/capability.svg}
\caption{Διαγράμματα ικανότητας σύγχρονης γεννήτριας (αριστερά) και μονάδας που συνδέεται μέσω
μετατροπέα (δεξιά), σε ανά μονάδα τιμές ως προς την ονομαστική φαινόμενη ισχύ και για τάση
1 p.u. Η περιοχή λειτουργίας της γεννήτριας είναι ασύμμετρη, με μεγαλύτερη ικανότητα
παραγωγής από ό,τι απορρόφησης άεργης ισχύος· η κουκκίδα είναι το ονομαστικό σημείο
λειτουργίας, με συντελεστή ισχύος 0,9. Η περιοχή του μετατροπέα είναι συμμετρική και
περιορίζεται από το ρεύμα του, οπότε συρρικνώνεται με την τάση· η διακεκομμένη καμπύλη
αντιστοιχεί σε τάση 0,9 p.u.}
\label{fig:capability}
\end{figure}

Τα υπόλοιπα μέσα ρύθμισης της τάσης είναι:

\begin{itemize}
\item \textbf{Σύγχρονοι αντισταθμιστές} (synchronous condensers). Είναι σύγχρονες μηχανές χωρίς
      κινητήρια μηχανή, που ρυθμίζουν την άεργη ισχύ τους μέσω της διέγερσης. Εκτός από άεργη
      ισχύ συνεισφέρουν ρεύμα βραχυκύκλωσης και αδράνεια, γι' αυτό η χρήση τους επεκτείνεται σε
      συστήματα με υψηλή διείσδυση μονάδων με μετατροπείς. Σε ορισμένες περιπτώσεις γεννήτριες
      σταθμών που αποσύρονται μετατρέπονται σε σύγχρονους αντισταθμιστές.
\item \textbf{Πυκνωτές και πηνία αντιστάθμισης.} Συνδέονται παράλληλα σε ζυγούς υποσταθμών ή
      στο τριτεύον τύλιγμα μετασχηματιστών και ζεύγνυνται μηχανικά, σε διακριτά βήματα. Οι
      πυκνωτές παράγουν άεργη ισχύ στις ώρες αιχμής και τα πηνία την απορροφούν στις ώρες
      χαμηλού φορτίου· τα πηνία των μεγάλων γραμμών συνδέονται συχνά μόνιμα στα άκρα τους. Η
      άεργη ισχύς ενός πυκνωτή, $Q = BV^2$, μειώνεται με το τετράγωνο της τάσης, δηλαδή ακριβώς
      όταν η στήριξη της τάσης είναι περισσότερο αναγκαία. Επιπλέον, ο αριθμός των χειρισμών
      είναι περιορισμένος, λόγω της καταπόνησης των διακοπτών και των μεταβατικών φαινομένων
      της ζεύξης.
\item \textbf{Πυκνωτές σειράς.} Αντισταθμίζουν μέρος της αντίδρασης σειράς μεγάλων γραμμών,
      οπότε μειώνουν τη διαφορά γωνιών και την άεργη ισχύ που καταναλώνει η γραμμή και
      αυξάνουν τη μεταφορική της ικανότητα. Η άεργη ισχύς που παράγουν, $X_C I^2$, αυξάνεται
      με το τετράγωνο του ρεύματος, δηλαδή με τη φόρτιση της γραμμής, οπότε η αντιστάθμιση
      προσαρμόζεται αυτόματα στις ανάγκες.
\item \textbf{Μετασχηματιστές με μεταγωγέα λήψεων υπό φορτίο} (On-Load Tap Changer, OLTC). Ο
      μεταγωγέας μεταβάλλει τον λόγο μετασχηματισμού σε βήματα της τάξης του 1\%, ώστε η τάση
      της μίας πλευράς να παραμένει εντός μιας ζώνης γύρω από την τιμή αναφοράς. Κάθε
      χειρισμός γίνεται με χρονική καθυστέρηση δεκάδων δευτερολέπτων, ώστε να αγνοούνται οι
      παροδικές μεταβολές. Ο μεταγωγέας δεν παράγει άεργη ισχύ, αλλά ανακατανέμει τις τάσεις
      και τις ροές άεργης ισχύος μεταξύ των δύο πλευρών. Η αύξηση της τάσης στην πλευρά του
      φορτίου αυξάνει τη ζήτηση, η οποία εξαρτάται από την τάση, και επομένως την άεργη ισχύ
      που απορροφάται από το δίκτυο μεταφοράς, με αποτέλεσμα να μειώνεται η τάση στην πλευρά
      του δικτύου μεταφοράς.
\item \textbf{Διατάξεις FACTS.} Οι SVC και STATCOM μεταβάλλουν την άεργη ισχύ τους συνεχώς και
      σε κλάσματα του δευτερολέπτου. Εξετάζονται στην Ενότητα «\nameref{sec:facts}».
\item \textbf{Χειρισμοί του δικτύου.} Σε συνθήκες πολύ χαμηλού φορτίου οι ΔΣΜ θέτουν ενίοτε
      εκτός λειτουργίας ελαφρά φορτισμένες γραμμές υπερυψηλής τάσης, ώστε να περιοριστεί η
      άεργη ισχύς που παράγουν. Ο χειρισμός αυτός μειώνει την εφεδρεία του δικτύου και
      επιτρέπεται μόνο εφόσον εξακολουθεί να ικανοποιείται το κριτήριο N-1 (Κεφάλαιο 1,
      Ενότητα «\nameref{sec:n1}»).
\end{itemize}

\subsection{Ευστάθεια τάσης}
\label{sec:v-stability}

Ευστάθεια τάσης (voltage stability) είναι η ικανότητα του συστήματος να διατηρεί αποδεκτές
τάσεις σε όλους τους ζυγούς μετά από μια διαταραχή, ξεκινώντας από δεδομένη κατάσταση
λειτουργίας. Η απώλειά της εκδηλώνεται συνήθως ως προοδευτική μείωση των τάσεων σε μια περιοχή
του συστήματος, η οποία μπορεί να καταλήξει σε κατάρρευση τάσης (voltage collapse) και σε
διακοπή της τροφοδότησης μεγάλου μέρους της ζήτησης. Εμφανίζεται κυρίως σε φορτισμένα
συστήματα, όπου η ισχύς μεταφέρεται σε μεγάλες αποστάσεις από τους σταθμούς παραγωγής προς τα
κέντρα κατανάλωσης, και οφείλεται στο ότι το δίκτυο δεν μπορεί να καλύψει τη ζήτηση ισχύος, και
ιδίως άεργης ισχύος, των φορτίων.

Το φαινόμενο εξηγείται με το σύστημα του Σχήματος~\ref{fig:vq} για $R = 0$. Ο ζυγός 1
παριστάνει το υπόλοιπο σύστημα ως πηγή σταθερής τάσης $E$, και ο ζυγός 2, με τάση $V$,
τροφοδοτεί φορτίο $P + jQ$ με σταθερό συντελεστή ισχύος, δηλαδή $Q = P\tan\varphi$. Το
τετράγωνο του μέτρου της \eqref{eq:v1phasor} δίνει $E^2 = (V + XQ/V)^2 + (XP/V)^2$ ή,
ισοδύναμα,

\begin{equation}
V^4 + \left(2QX - E^2\right) V^2 + X^2\left(P^2 + Q^2\right) = 0
\label{eq:pv}
\end{equation}

Η \eqref{eq:pv} είναι δευτεροβάθμια ως προς $V^2$. Για κάθε τιμή της $P$ μέχρι ένα μέγιστο
έχει δύο λύσεις, δηλαδή δύο σημεία λειτουργίας: ένα με υψηλή τάση και μικρό ρεύμα, στο οποίο
λειτουργεί κανονικά το σύστημα, και ένα με χαμηλή τάση και μεγάλο ρεύμα. Καθώς η $P$ αυξάνεται,
οι δύο λύσεις πλησιάζουν και συμπίπτουν όταν μηδενιστεί η διακρίνουσα, στη μέγιστη ισχύ που
μπορεί να μεταφερθεί:

\begin{equation}
P_{\max} = \frac{E^2 \cos\varphi}{2X\left(1 + \sin\varphi\right)}
\label{eq:pmax}
\end{equation}

Για μοναδιαίο συντελεστή ισχύος $P_{\max} = E^2/(2X)$ και η αντίστοιχη τάση είναι
$E/\sqrt{2}$, δηλαδή περίπου 0,71 της τάσης της πηγής.

Η γραφική παράσταση της $V$ ως συνάρτησης της $P$ ονομάζεται καμπύλη $P$--$V$ ή, λόγω του
σχήματός της, καμπύλη μύτης (nose curve). Το Σχήμα~\ref{fig:pv} δείχνει τις καμπύλες σε
κανονικοποιημένη μορφή για διάφορους συντελεστές ισχύος. Το επάνω τμήμα κάθε καμπύλης είναι η
περιοχή κανονικής λειτουργίας. Στο κάτω τμήμα η λειτουργία είναι ασταθής για φορτία που τείνουν
να αποκαταστήσουν την ισχύ τους, όπως εξηγείται παρακάτω. Η απόσταση του σημείου λειτουργίας
από τη μύτη της καμπύλης, το \emph{περιθώριο φόρτισης} (loadability margin), είναι το βασικό
μέτρο της εγγύτητας στην αστάθεια τάσης.

\begin{figure}[!htbp]
\centering
\includegraphics[width=\linewidth]{figures/pv-curves.svg}
\caption{Καμπύλες $P$--$V$ του συστήματος δύο ζυγών για διάφορες τιμές του $\tan\varphi$ του
φορτίου, σε κανονικοποιημένη μορφή. Θετικές τιμές αντιστοιχούν σε επαγωγικό φορτίο και
αρνητικές σε φορτίο με αντιστάθμιση μεγαλύτερη από την άεργη ζήτησή του. Οι συνεχείς γραμμές
είναι τα επάνω τμήματα, στα οποία λειτουργεί κανονικά το σύστημα, και οι διακεκομμένες τα
κάτω. Οι κουκκίδες είναι τα κρίσιμα σημεία μέγιστης μεταφερόμενης ισχύος.}
\label{fig:pv}
\end{figure}

Από το σχήμα προκύπτουν δύο συμπεράσματα. Πρώτον, τα επαγωγικά φορτία μειώνουν τη μέγιστη
ισχύ, ενώ η αντιστάθμιση της άεργης ζήτησής τους την αυξάνει. Δεύτερον, με την αντιστάθμιση
η τάση στη μύτη της καμπύλης, η \emph{κρίσιμη τάση}, αυξάνεται και πλησιάζει τις κανονικές
τιμές λειτουργίας. Σε ισχυρά αντισταθμισμένο σύστημα η τάση μπορεί να παραμένει αποδεκτή μέχρι
πολύ κοντά στο όριο, οπότε το μέτρο της τάσης δεν αποτελεί από μόνο του αξιόπιστο δείκτη της
εγγύτητας στην κατάρρευση.

Συμπληρωματική εικόνα δίνει η καμπύλη $Q$--$V$ ενός ζυγού: η άεργη ισχύς $Q_c$ που πρέπει να
εγχέει ένας πλασματικός αντισταθμιστής στον ζυγό, ώστε η τάση του να είναι $V$, για δεδομένη
ζήτηση $P + jQ$. Για το σύστημα δύο ζυγών, με απαλοιφή της $\delta$ από την
\eqref{eq:pq-line},

\begin{equation}
Q_c = Q + \frac{V^2}{X} - \sqrt{\left(\frac{EV}{X}\right)^2 - P^2}
\label{eq:qv}
\end{equation}

Οι καμπύλες του Σχήματος~\ref{fig:qv} παρουσιάζουν ελάχιστο. Δεξιά του ελαχίστου η αύξηση της
εγχεόμενης άεργης ισχύος αυξάνει την τάση ($\mathrm{d}Q_c/\mathrm{d}V > 0$), όπως στην
κανονική λειτουργία. Αριστερά του ελαχίστου η σχέση αντιστρέφεται, οπότε η ρύθμιση της τάσης
με έγχυση άεργης ισχύος οδηγεί σε αστάθεια. Το σημείο λειτουργίας χωρίς αντιστάθμιση είναι η
τομή του δεξιού κλάδου με τον άξονα $Q_c = 0$, και η απόσταση του ελαχίστου από τον άξονα
αυτόν είναι το \emph{περιθώριο άεργης ισχύος} του ζυγού: η επιπλέον άεργη ζήτηση που μπορεί να
καλυφθεί πριν από την αστάθεια. Όταν το ελάχιστο βρίσκεται πάνω από τον άξονα, το σύστημα δεν
μπορεί να λειτουργήσει χωρίς αντιστάθμιση. Στην πράξη οι καμπύλες $Q$--$V$ υπολογίζονται με
διαδοχικές επιλύσεις ροής ισχύος, στις οποίες ο εξεταζόμενος ζυγός θεωρείται ζυγός ελεγχόμενης
τάσης χωρίς όρια άεργης ισχύος.

\begin{figure}[!htbp]
\centering
\includegraphics[width=\linewidth]{figures/qv-curves.svg}
\caption{Καμπύλες $Q$--$V$ στον ζυγό φορτίου του συστήματος δύο ζυγών, σε κανονικοποιημένη
μορφή, για φορτίο με $\tan\varphi$ = 0,2 και διάφορες τιμές της ενεργής ισχύος. Οι κουκκίδες
είναι τα ελάχιστα, που χωρίζουν τον ευσταθή δεξιό κλάδο από τον ασταθή αριστερό
(διακεκομμένη γραμμή), και οι κύκλοι τα σημεία λειτουργίας χωρίς αντιστάθμιση.}
\label{fig:qv}
\end{figure}

Η αστάθεια τάσης οφείλεται κυρίως στη δυναμική των φορτίων. Αμέσως μετά από μια διαταραχή, για
παράδειγμα την απώλεια μιας γραμμής, η καμπύλη $P$--$V$ συρρικνώνεται, η τάση μειώνεται και
μαζί της μειώνεται η ζήτηση, η οποία εξαρτάται από την τάση. Στη συνέχεια διάφοροι μηχανισμοί
τείνουν να αποκαταστήσουν την ισχύ των φορτίων: οι μεταγωγείς λήψεων των μετασχηματιστών
διανομής επαναφέρουν σε μερικά λεπτά την τάση στην πλευρά του φορτίου, τα θερμοστατικά
ελεγχόμενα φορτία παραμένουν περισσότερο σε λειτουργία και οι επαγωγικοί κινητήρες αυξάνουν
την ολίσθησή τους. Αν η ζήτηση που τείνει να αποκατασταθεί υπερβαίνει τη μέγιστη ισχύ της νέας
καμπύλης, δεν υπάρχει σημείο ισορροπίας και οι τάσεις μειώνονται προοδευτικά. Η πορεία
επιταχύνεται όταν οι γεννήτριες φθάνουν στο όριο διέγερσης και ο OEL περιορίζει την άεργη ισχύ
τους. Ανάλογα με τη χρονική κλίμακα των μηχανισμών, η αστάθεια χαρακτηρίζεται
\emph{βραχυπρόθεσμη}, με διάρκεια λίγων δευτερολέπτων, όταν οφείλεται σε επαγωγικούς κινητήρες
ή σε μετατροπείς, και \emph{μακροπρόθεσμη}, με διάρκεια λεπτών, όταν οφείλεται σε μεταγωγείς
λήψεων, θερμοστατικά φορτία και περιοριστές υπερδιέγερσης.

Τα μέτρα πρόληψης περιλαμβάνουν τη διατήρηση επαρκούς εφεδρείας άεργης ισχύος, ιδίως ταχείας,
κοντά στα κέντρα κατανάλωσης, τη ζεύξη πυκνωτών και τη συντονισμένη ρύθμιση της τάσης της
επόμενης ενότητας. Σε καταστάσεις έκτακτης ανάγκης εφαρμόζονται η αναστολή των χειρισμών των
μεταγωγέων λήψεων και η απόρριψη φορτίου λόγω υπότασης (undervoltage load shedding).

\subsection{Ιεραρχική ρύθμιση τάσης}
\label{sec:v-hier}

Η ρύθμιση της τάσης οργανώνεται ιεραρχικά σε τρία επίπεδα, τα οποία διαφέρουν ως προς τη
γεωγραφική έκταση και τη χρονική κλίμακα (Πίνακας~\ref{tab:v-hier}). Κάθε επίπεδο είναι
σημαντικά βραδύτερο από το προηγούμενο, ώστε να αντιλαμβάνεται τα ταχύτερα επίπεδα σε μόνιμη
κατάσταση και να αποφεύγονται ανεπιθύμητες αλληλεπιδράσεις.

\begin{itemize}
\item \textbf{Πρωτεύουσα ρύθμιση.} Τοπική και αυτόματη, από τους AVR των γεννητριών, τις
      διατάξεις FACTS και τους μετατροπείς με ρύθμιση τάσης. Αποκρίνεται σε κλάσματα του
      δευτερολέπτου έως λίγα δευτερόλεπτα και διατηρεί την τάση κάθε ρυθμιζόμενου ζυγού κοντά
      στην τιμή αναφοράς του.
\item \textbf{Δευτερεύουσα ρύθμιση.} Περιφερειακή, σε χρονική κλίμακα λεπτών. Το σύστημα
      χωρίζεται σε ζώνες και η τάση ενός αντιπροσωπευτικού ζυγού κάθε ζώνης, του
      \emph{ζυγού-πιλότου} (pilot bus), ρυθμίζεται με μεταβολή των τιμών αναφοράς των AVR των
      γεννητριών της ζώνης, ώστε όλες να συμμετέχουν με το ίδιο ποσοστό της ικανότητάς τους.
      Αυτόματη δευτερεύουσα ρύθμιση εφαρμόζεται σε ορισμένα συστήματα, όπως της Γαλλίας και
      της Ιταλίας. Σε πολλά συστήματα τον ρόλο αυτό εκτελούν οι χειριστές του κέντρου ελέγχου,
      με ζεύξη πυκνωτών και πηνίων, αλλαγή των τιμών αναφοράς των γεννητριών και χειρισμούς
      των μεταγωγέων λήψεων.
\item \textbf{Τριτεύουσα ρύθμιση.} Κεντρική, σε χρονική κλίμακα δεκάδων λεπτών έως ωρών. Με
      βέλτιστη ροή ισχύος (Optimal Power Flow, OPF) υπολογίζονται οι τιμές αναφοράς των
      τάσεων που ελαχιστοποιούν τις απώλειες και διατηρούν επαρκή εφεδρεία άεργης ισχύος,
      τηρώντας τα όρια λειτουργίας και το κριτήριο N-1.
\end{itemize}

\begin{table}[!htbp]
\centering
\begin{tabular}{llll}
\hline
Επίπεδο & Έκταση & Χρονική κλίμακα & Κύρια μέσα \\
\hline
πρωτεύουσα   & τοπική  & ms -- s     & AVR, SVC, STATCOM, μετατροπείς \\
δευτερεύουσα & ζώνη    & min         & τιμές αναφοράς AVR, πυκνωτές, πηνία \\
τριτεύουσα   & σύστημα & 15 min -- h & βέλτιστη ροή ισχύος \\
\hline
\end{tabular}
\caption{Επίπεδα της ιεραρχικής ρύθμισης τάσης.}
\label{tab:v-hier}
\end{table}

Δύο αρχές συντονισμού είναι ιδιαίτερα σημαντικές. Πρώτον, η ταχεία εφεδρεία άεργης ισχύος των
γεννητριών και των διατάξεων FACTS πρέπει να διατηρείται για την αντιμετώπιση διαταραχών. Τα
βραδέα μέσα, όπως οι πυκνωτές και τα πηνία, αναλαμβάνουν σταδιακά τη μόνιμη κάλυψη των αναγκών,
ώστε οι ταχείες πηγές να επανέρχονται κοντά στο μέσο της περιοχής λειτουργίας τους. Δεύτερον,
σε μετασχηματιστές σε σειρά, όπως ένας μετασχηματιστής 150/20 kV που τροφοδοτείται από
μετασχηματιστή 400/150 kV, ο μεταγωγέας του ανώτερου επιπέδου τάσης πρέπει να δρα πρώτος, με
μικρότερη χρονική καθυστέρηση, ώστε να αποφεύγονται περιττοί χειρισμοί στο κατώτερο επίπεδο. Ο
συντονισμός της άεργης ισχύος στα σημεία σύνδεσης του συστήματος μεταφοράς με τα δίκτυα
διανομής αφορά και τους δύο διαχειριστές και εξετάζεται στο Κεφάλαιο 7.

\section{Ευέλικτα συστήματα μεταφοράς εναλλασσόμενου ρεύματος (FACTS)}
\label{sec:facts}

Πέρα από τη ρύθμιση της παραγωγής μέσω των αγορών και των εφεδρειών συχνότητας, ο Διαχειριστής
Συστήματος Μεταφοράς (ΔΣΜ) ελέγχει άμεσα τις ροές ισχύος και τις τάσεις του δικτύου μεταφοράς. Οι διατάξεις
\textbf{FACTS} είναι διατάξεις ηλεκτρονικών ισχύος που
επιτρέπουν τον ταχύ και συνεχή έλεγχο των μεγεθών που καθορίζουν τη ροή ισχύος σε δίκτυο
εναλλασσόμενου ρεύματος.

Σύμφωνα με την \eqref{eq:pq-line}, η ενεργή ισχύς που μεταφέρεται μεταξύ δύο ζυγών μέσω
γραμμής χωρίς απώλειες εξαρτάται από τα μέτρα των τάσεων $V_1, V_2$, την αντίδραση σειράς
$X$ και τη διαφορά γωνιών $\theta_1 - \theta_2$. Σε
συμβατικό δίκτυο τα μεγέθη αυτά δεν ελέγχονται άμεσα, οπότε η κατανομή των ροών καθορίζεται
από τις σύνθετες αντιστάσεις των κλάδων και όχι από τις απαιτήσεις λειτουργίας. Οι διατάξεις
FACTS καθιστούν ελεγχόμενο ένα ή περισσότερα από τα μεγέθη αυτά και κατατάσσονται ανάλογα με
το μέγεθος που επηρεάζουν (Σχήμα~\ref{fig:facts}):

\begin{itemize}
\item \textbf{Παράλληλης σύνδεσης} (έλεγχος του $V$): ο στατός αντισταθμιστής άεργης ισχύος
      (Static Var Compensator, SVC) και ο στατός σύγχρονος αντισταθμιστής (Static
      Synchronous Compensator, STATCOM). Ρυθμίζουν την τάση του ζυγού μέσω έγχυσης ή
      απορρόφησης άεργης ισχύος.
\item \textbf{Σειριακής σύνδεσης} (έλεγχος του $X$): ο πυκνωτής σειράς ελεγχόμενος με
      θυρίστορ (Thyristor Controlled Series Capacitor, TCSC) και ο στατός σύγχρονος
      αντισταθμιστής σειράς (Static Synchronous Series Compensator, SSSC). Μεταβάλλουν την
      ισοδύναμη αντίδραση σειράς της γραμμής, ανακατανέμουν τις ροές ισχύος και αυξάνουν τη
      μεταφορική ικανότητα.
\item \textbf{Συνδυασμένης σύνδεσης}, σε σειρά και παράλληλα (έλεγχος και της γωνίας): ο
      ενοποιημένος ελεγκτής ροής ισχύος (Unified Power Flow Controller, UPFC) ελέγχει
      ταυτόχρονα και ανεξάρτητα την τάση και τις ροές ενεργής και άεργης ισχύος. Παρέχει τη
      μεγαλύτερη ευελιξία ελέγχου, με το υψηλότερο κόστος.
\end{itemize}

\begin{figure}[!htbp]
\centering
\includegraphics[width=\linewidth]{figures/facts.svg}
\caption{Οι τρεις κατηγορίες διατάξεων FACTS σε γραμμή δύο ζυγών. Κάθε κατηγορία επεμβαίνει
σε διαφορετικό μέγεθος της σχέσης μεταφοράς ισχύος, το οποίο σημειώνεται με έντονο χρώμα: η
παράλληλη στην τάση, η σειριακή στην αντίδραση, η συνδυασμένη και στη γωνία.}
\label{fig:facts}
\end{figure}

\subsection{Τοπολογίες διατάξεων}

Ανάλογα με την αρχή λειτουργίας τους, οι διατάξεις FACTS διακρίνονται σε διατάξεις
\emph{μεταβλητής σύνθετης αντίστασης} (εμπέδησης), με θυρίστορ, και σε διατάξεις
\emph{ελεγχόμενης πηγής τάσης}, με μετατροπείς πηγής τάσης. Τα μονογραμμικά διαγράμματα
δίνονται στο Σχήμα~\ref{fig:facts-devices}.

\begin{itemize}
\item Ο \textbf{SVC} αποτελείται από πηνίο ελεγχόμενο με θυρίστορ (Thyristor Controlled
      Reactor, TCR) και πυκνωτές μεταγόμενους με θυρίστορ (Thyristor Switched Capacitor,
      TSC), σε παράλληλη σύνδεση. Η γωνία έναυσης του TCR ρυθμίζει συνεχώς την ισοδύναμη
      επαγωγική επιδεκτικότητα, ενώ οι TSC προσθέτουν χωρητική επιδεκτικότητα σε διακριτά
      βήματα. Η διάταξη συμπεριφέρεται ως \emph{μεταβλητή επιδεκτικότητα} $B$ (susceptance).
\item Ο \textbf{STATCOM} είναι μετατροπέας πηγής τάσης (Voltage Source Converter, VSC) με
      πυκνωτή στην πλευρά συνεχούς ρεύματος (ΣΡ). Παράγει εναλλασσόμενη τάση $V_\mu$
      ελεγχόμενου μέτρου, η οποία εφαρμόζεται στο δίκτυο μέσω της αντίδρασης σύζευξης: για
      $V_\mu > V$ ο STATCOM εγχέει άεργη ισχύ, ενώ για $V_\mu < V$ απορροφά.
\item Ο \textbf{TCSC} είναι πυκνωτής σειράς με παράλληλο TCR, οπότε η ισοδύναμη αντίδραση
      σειράς της γραμμής μεταβάλλεται συνεχώς. Ο \textbf{SSSC} εγχέει τάση σε σειρά μέσω
      μετασχηματιστή σύζευξης σειράς· η εγχεόμενη τάση είναι ανεξάρτητη του ρεύματος της
      γραμμής, ενώ στον TCSC η τάση στα άκρα του πυκνωτή είναι ανάλογη του ρεύματος.
\item Ο \textbf{UPFC} αποτελείται από έναν παράλληλο και έναν σειριακό VSC με \emph{κοινό
      ζυγό ΣΡ} (DC link). Ο κοινός ζυγός επιτρέπει την ανταλλαγή \emph{ενεργής} ισχύος
      μεταξύ των δύο μετατροπέων, οπότε ο σειριακός μετατροπέας εγχέει τάση οποιουδήποτε
      μέτρου και γωνίας εντός των ονομαστικών του ορίων και ελέγχει ανεξάρτητα τις ροές
      ενεργής και άεργης ισχύος.
\end{itemize}

\begin{figure}[!htbp]
\centering
\includegraphics[width=\linewidth]{figures/facts-devices.svg}
\caption{Μονογραμμικά διαγράμματα των βασικών διατάξεων FACTS. Οι δύο πρώτες συνδέονται
παράλληλα, οι δύο επόμενες σε σειρά, ενώ ο UPFC συνδυάζει αμφότερους τους τρόπους μέσω
κοινού ζυγού ΣΡ.}
\label{fig:facts-devices}
\end{figure}

\subsection{Χαρακτηριστικές V--I των SVC και STATCOM}

Η διαφορά στην αρχή λειτουργίας αποτυπώνεται στις χαρακτηριστικές $V$--$I$ των δύο
διατάξεων παράλληλης σύνδεσης (Σχήμα~\ref{fig:facts-vi}).

Ο SVC συμπεριφέρεται ως μεταβλητή επιδεκτικότητα, οπότε $I = B\,V$ και $Q = B\,V^2$. Στο
όριο χωρητικής λειτουργίας το ρεύμα είναι ανάλογο της τάσης του ζυγού: στο 50\% της
ονομαστικής τάσης ο SVC αποδίδει το 50\% του ονομαστικού ρεύματος και το 25\% της ονομαστικής
άεργης ισχύος. Ο STATCOM, ως πηγή τάσης με έλεγχο ρεύματος, διατηρεί το ονομαστικό ρεύμα έως
πολύ χαμηλές τιμές τάσης, με περιορισμό μόνο από το μέγιστο επιτρεπόμενο ρεύμα των
ημιαγωγικών διακοπτών· η άεργη ισχύς του μειώνεται επομένως γραμμικά και όχι τετραγωνικά
με την τάση.

Κατά τη διάρκεια και αμέσως μετά από βραχυκύκλωμα, όταν η τάση είναι χαμηλή και η ανάγκη
στήριξης τάσης μέγιστη, ο STATCOM αποδίδει σημαντικά μεγαλύτερη άεργη ισχύ από SVC ίδιας
ονομαστικής ισχύος.

\begin{figure}[!htbp]
\centering
\includegraphics[width=\linewidth]{figures/facts-vi.svg}
\caption{Χαρακτηριστικά $V$--$I$ των δύο διατάξεων παράλληλης σύνδεσης. Η σκιασμένη περιοχή
είναι η περιοχή λειτουργίας και η έντονη γραμμή η χαρακτηριστική λειτουργίας με στατισμό. Τα
όρια του SVC είναι ευθείες σταθερής επιδεκτικότητας που διέρχονται από την αρχή, ενώ του
STATCOM είναι κατακόρυφες ευθείες σταθερού ρεύματος.}
\label{fig:facts-vi}
\end{figure}

Οι διατάξεις FACTS, μαζί με τα συστήματα συνεχούς ρεύματος υψηλής τάσης (High Voltage
Direct Current, HVDC), στα οποία η ροή ισχύος καθορίζεται απευθείας από τον έλεγχο των
μετατροπέων, συμπληρώνουν τα συμβατικά μέσα του συστήματος μεταφοράς: για τη ρύθμιση της τάσης
τους ρυθμιστές τάσης των γεννητριών, τα πηνία και τους πυκνωτές αντιστάθμισης και τους
μεταγωγείς λήψεων των μετασχηματιστών, και για τη διαχείριση των συμφορήσεων την ανακατανομή
παραγωγής, τις αλλαγές τοπολογίας και τους μετασχηματιστές μετατόπισης φάσης. Τα αντίστοιχα μέσα στα
δίκτυα διανομής εξετάζονται στα Κεφάλαια 5 και 6.

\section{Γρήγορη έγχυση άεργου ρεύματος}
\label{sec:frt}

Κατά τη διάρκεια ενός βραχυκυκλώματος η τάση μειώνεται απότομα σε μια περιοχή του συστήματος,
η έκταση της οποίας εξαρτάται από την ισχύ βραχυκύκλωσης του δικτύου. Η κύρια προστασία
απομονώνει το σφάλμα σε χρόνο της τάξης των 100 ms, αλλά οι επιπτώσεις του μπορεί να είναι
πολύ μεγαλύτερες. Αν οι μονάδες ηλεκτροπαραγωγής της περιοχής αποσυνδεθούν λόγω της βύθισης
της τάσης, ένα τοπικό σφάλμα στο δίκτυο μεταφοράς μετατρέπεται σε απώλεια παραγωγής, η οποία
μπορεί να υπερβεί το συμβάν αναφοράς για το οποίο διαστασιολογούνται οι εφεδρείες συχνότητας
(\hyperref[kef-03-theoria]{Κεφάλαιο 3}). Η υπηρεσία έχει επομένως δύο σκέλη: οι μονάδες πρέπει
να παραμένουν συνδεδεμένες κατά τη διάρκεια του σφάλματος και να εγχέουν άεργο ρεύμα που
στηρίζει την τάση. Η στήριξη αυτή περιορίζει την έκταση της βύθισης, βοηθά την προστασία να
ανιχνεύσει το σφάλμα και επιταχύνει την ανάκτηση της τάσης μετά την απομόνωσή του.

\subsection{Απόκριση των σύγχρονων μηχανών}
\label{sec:sm-response}

Η σύγχρονη μηχανή παρέχει την υπηρεσία εγγενώς. Αμέσως μετά από μια μεταβολή της τάσης στους
ακροδέκτες της, συμπεριφέρεται ως πηγή τάσης πίσω από την υπομεταβατική αντίδραση $X_d^{\prime\prime}$,
τυπικά 0,15--0,25 p.u. Μια βύθιση της τάσης κατά $\Delta V$ προκαλεί επομένως ακαριαία
μεταβολή του ρεύματος

\begin{equation}
\Delta I \approx \frac{\Delta V}{X_d^{\prime\prime}}
\label{eq:sm-dI}
\end{equation}

σε ανά μονάδα τιμές, χωρίς καμία ενέργεια ελέγχου. Σε τριφασικό βραχυκύκλωμα στους ακροδέκτες
το αρχικό ρεύμα φθάνει τις 4 έως 7 φορές το ονομαστικό και στη συνέχεια μειώνεται, καθώς η
μηχανή περνά από την υπομεταβατική στη μεταβατική και στη μόνιμη κατάσταση. Στο μεταξύ ο
ρυθμιστής τάσης αυξάνει τη διέγερση στη μέγιστη τιμή της (field forcing) και συγκρατεί την
τάση. Επειδή η σύνθετη αντίσταση του δικτύου και του σφάλματος είναι κυρίως επαγωγική, το ρεύμα
αυτό είναι κατά κύριο λόγο άεργο.

\subsection{Παραμονή σε λειτουργία κατά τα σφάλματα}
\label{sec:frt-req}

Οι μονάδες που συνδέονται μέσω μετατροπέων συμπεριφέρονται διαφορετικά. Το ρεύμα του
μετατροπέα μπορεί να υπερβεί το ονομαστικό μόνο κατά 10--20\%, γιατί οι ημιαγωγικοί διακόπτες
έχουν πολύ μικρή θερμική χωρητικότητα και δεν αντέχουν υπερφόρτιση ούτε για λίγα ms. Επιπλέον,
το ρεύμα καθορίζεται από τον έλεγχο: ένας μετατροπέας που ακολουθεί το δίκτυο (grid-following)
μετρά τη γωνία της τάσης με βρόχο κλειδωμένης φάσης (Phase-Locked Loop, PLL) και εγχέει το
ρεύμα που ορίζουν οι ελεγκτές του. Οι πρώτες ανεμογεννήτριες αποσυνδέονταν σε κάθε βύθιση της
τάσης, για να προστατευθούν. Με την αύξηση της αιολικής παραγωγής, η μαζική αποσύνδεση μετά
από ένα σφάλμα στο δίκτυο μεταφοράς έγινε απειλή για την ευστάθεια του συστήματος, και οι
κώδικες σύνδεσης εισήγαγαν την απαίτηση παραμονής σε λειτουργία κατά τα σφάλματα (Fault Ride
Through, FRT).

\begin{figure}[!htbp]
\centering
\includegraphics[width=\linewidth]{figures/frt.svg}
\caption{Αριστερά: ενδεικτική καμπύλη παραμονής σε λειτουργία για μονάδα που συνδέεται μέσω
μετατροπέα και μια τυπική τροχιά της τάσης, η οποία πέφτει στα 0,25 p.u. κατά το σφάλμα και
ανακάμπτει μετά την απομόνωσή του στα 120 ms. Δεξιά: χαρακτηριστική έγχυσης άεργου ρεύματος
μετατροπέα με $k$ = 2, η οποία περιορίζεται από το μέγιστο ρεύμα του, και εγγενής απόκριση
σύγχρονης μηχανής με $X_d^{\prime\prime}$ = 0,2 p.u.}
\label{fig:frt}
\end{figure}

Η απαίτηση διατυπώνεται ως καμπύλη τάσης--χρόνου στο σημείο σύνδεσης (Σχήμα~\ref{fig:frt},
αριστερά). Η μονάδα πρέπει να παραμένει συνδεδεμένη και σε ευσταθή λειτουργία όσο η τάση
βρίσκεται πάνω από την καμπύλη, ενώ επιτρέπεται να αποσυνδεθεί όταν η τάση πέσει κάτω από
αυτήν. Ο RfG κατατάσσει τις μονάδες σε τέσσερις τύπους, A έως D, ανάλογα με την ισχύ τους και
την τάση σύνδεσης· στον τύπο D ανήκουν οι μονάδες που συνδέονται σε τάση 110 kV και άνω. Για
κάθε τύπο ορίζει τη μορφή της καμπύλης και τα εύρη των παραμέτρων της, ενώ τις τιμές καθορίζει
ο ΔΣΜ. Για μονάδες τύπου D που συνδέονται μέσω μετατροπέων, η καμπύλη μπορεί να απαιτεί
παραμονή σε λειτουργία ακόμη και με μηδενική τάση για 140--150 ms, χρόνο που καλύπτει την
απομόνωση του σφάλματος από την κύρια προστασία, και στη συνέχεια όσο η τάση ανακάμπτει πάνω
από μια ευθεία που φθάνει τα 0,85 p.u. σε 1,5--3 s. Πολλοί κώδικες ορίζουν αντίστοιχες
απαιτήσεις και για τις υπερτάσεις που ακολουθούν την απομόνωση του σφάλματος.

\subsection{Έγχυση άεργου ρεύματος από μετατροπείς}
\label{sec:kfactor}

Πέρα από την παραμονή σε λειτουργία, οι κώδικες σύνδεσης απαιτούν από τις μονάδες με
μετατροπείς να στηρίζουν την τάση κατά τη διάρκεια της βύθισης με άεργο ρεύμα ανάλογο της
απόκλισης της τάσης:

\begin{equation}
\Delta I_q = k\,\Delta U
\label{eq:kfactor}
\end{equation}

όπου $\Delta U = U_0 - U$ η βύθιση της τάσης ως προς την τιμή $U_0$ πριν από το σφάλμα, σε ανά
μονάδα της ονομαστικής τάσης, $\Delta I_q$ η αύξηση του χωρητικού άεργου ρεύματος, σε ανά
μονάδα του ονομαστικού ρεύματος, και $k$ ο συντελεστής ενίσχυσης, με τυπική τιμή 2: κάθε 1\%
βύθισης της τάσης απαιτεί 2\% επιπλέον άεργο ρεύμα. Σε ορισμένους κώδικες η χαρακτηριστική
έχει νεκρή ζώνη γύρω από την ονομαστική τάση, ώστε η στήριξη να ενεργοποιείται μόνο σε
πραγματικές βυθίσεις. Η απόκριση πρέπει να ολοκληρώνεται σε μερικές δεκάδες ms. Σε υπερτάσεις,
όπου $\Delta U < 0$, η ίδια χαρακτηριστική δίνει επαγωγικό ρεύμα. Η σύγκριση με την
\eqref{eq:sm-dI} δείχνει ότι η σύγχρονη μηχανή αποκρίνεται με ισοδύναμο συντελεστή $1/X_d^{\prime\prime}$,
δηλαδή 4 έως 7, και χωρίς το όριο ρεύματος του μετατροπέα (Σχήμα~\ref{fig:frt}, δεξιά).

Επειδή το συνολικό ρεύμα του μετατροπέα είναι περιορισμένο,

\begin{equation}
I_p^2 + I_q^2 \le I_{\max}^2
\label{eq:ilimit}
\end{equation}

όπου $I_p$ και $I_q$ η ενεργή και η άεργη συνιστώσα του ρεύματος, κατά τη διάρκεια της βύθισης
δίνεται προτεραιότητα στο άεργο ρεύμα και το ενεργό ρεύμα μειώνεται όσο χρειάζεται. Η ενεργή
ισχύς μειώνεται επομένως για δύο λόγους, επειδή μειώνονται τόσο η τάση όσο και το ενεργό
ρεύμα. Η ισχύς της πρωτογενούς πηγής που δεν εγχέεται στο δίκτυο αποθηκεύεται προσωρινά ως
κινητική ενέργεια στον δρομέα της ανεμογεννήτριας, καταναλώνεται σε αντιστάτη πέδησης στον ζυγό
ΣΡ ή, στα φωτοβολταϊκά, περιορίζεται με μετακίνηση του σημείου λειτουργίας των πλαισίων. Μετά
την απομόνωση του σφάλματος η ενεργή ισχύς πρέπει να αποκατασταθεί σε χρόνο που ορίζει ο ΔΣΜ.
Η γρήγορη αποκατάσταση περιορίζει τη διαταραχή της συχνότητας, αλλά σε ασθενή δίκτυα
επιβαρύνει την ανάκτηση της τάσης.

\emph{Παράδειγμα.} Αιολικό πάρκο λειτουργεί με ονομαστική ενεργή ισχύ και μοναδιαίο συντελεστή
ισχύος, όταν ένα σφάλμα μειώνει την τάση του σημείου σύνδεσης στα 0,5 p.u. Με $k$ = 2 και
χωρίς νεκρή ζώνη, η \eqref{eq:kfactor} απαιτεί άεργο ρεύμα 2 × 0,5 = 1,0 p.u. Αν το μέγιστο
ρεύμα του μετατροπέα είναι 1,1 p.u., η \eqref{eq:ilimit} περιορίζει το ενεργό ρεύμα σε 0,46
p.u. περίπου και την ενεργή ισχύ σε 0,5 × 0,46 ≈ 0,23 p.u. Το πάρκο εγχέει άεργη ισχύ
0,5 × 1,0 = 0,5 p.u., δηλαδή το μισό από ό,τι θα έδινε με το ίδιο ρεύμα σε ονομαστική τάση.

Τα περισσότερα σφάλματα είναι μονοφασικά και προκαλούν ασύμμετρες βυθίσεις. Ένας μετατροπέας
που εγχέει μόνο ρεύμα θετικής ακολουθίας δεν περιορίζει την τάση αρνητικής ακολουθίας και
μπορεί να δυσχεράνει τη λειτουργία της προστασίας, η οποία χρησιμοποιεί συχνά μεγέθη αρνητικής
ακολουθίας. Γι' αυτό νεότεροι κώδικες απαιτούν και έγχυση ρεύματος αρνητικής ακολουθίας, κατ'
αναλογία με την \eqref{eq:kfactor}.

Ως κύρια λειτουργία τους, τη γρήγορη έγχυση άεργου ρεύματος παρέχουν οι STATCOM και οι
σύγχρονοι αντισταθμιστές. Ο STATCOM διατηρεί το ονομαστικό του ρεύμα σε πολύ χαμηλές τάσεις
(Ενότητα «\nameref{sec:facts}»), ενώ ο σύγχρονος αντισταθμιστής δίνει στους πρώτους κύκλους,
σύμφωνα με την \eqref{eq:sm-dI}, ρεύμα πολλαπλάσιο του ονομαστικού και συμβάλλει επιπλέον στο
ρεύμα βραχυκύκλωσης και στην αδράνεια. Η ταχεία στήριξη βοηθά και την ανάκτηση της τάσης μετά
το σφάλμα, όταν οι επαγωγικοί κινητήρες που επιβραδύνθηκαν απορροφούν μεγάλο άεργο ρεύμα για να
επιταχυνθούν, φαινόμενο που συνδέεται με τη βραχυπρόθεσμη αστάθεια τάσης (Ενότητα
«\nameref{sec:v-stability}»).

\section{Ρεύμα βραχυκύκλωσης και ισχύς του δικτύου}
\label{sec:ssc}

Η ισχύς βραχυκύκλωσης ενός ζυγού

\begin{equation}
S_{\mathrm{sc}} = \sqrt{3}\, U_n I_{\mathrm{sc}}
\label{eq:ssc}
\end{equation}

όπου $U_n$ η ονομαστική πολική τάση και $I_{\mathrm{sc}}$ το ρεύμα τριφασικού βραχυκυκλώματος
στον ζυγό, εκφράζει πόσο «ισχυρό» είναι το δίκτυο στο σημείο αυτό. Σε ανά μονάδα τιμές είναι
το αντίστροφο της σύνθετης αντίστασης Thévenin του δικτύου, όπως φαίνεται από τον ζυγό, και,
σύμφωνα με την \eqref{eq:dv-ssc}, καθορίζει την ευαισθησία της τάσης στις μεταβολές της άεργης
ισχύος. Το ρεύμα βραχυκύκλωσης προέρχεται κυρίως από τις σύγχρονες μηχανές που βρίσκονται
ηλεκτρικά κοντά στον ζυγό, οι οποίες δίνουν 4 έως 7 φορές το ονομαστικό τους ρεύμα, ενώ οι
μονάδες με μετατροπείς συνεισφέρουν 1 έως 1,2 φορές. Καθώς οι σύγχρονες γεννήτριες
αντικαθίστανται από μονάδες με μετατροπείς, η ισχύς βραχυκύκλωσης μειώνεται, ιδίως σε περιοχές
απομακρυσμένες από τους μεγάλους σταθμούς και στις ώρες υψηλής ανανεώσιμης παραγωγής.

Επαρκής ισχύς βραχυκύκλωσης απαιτείται για:

\begin{itemize}
\item \textbf{την προστασία}, γιατί οι ηλεκτρονόμοι υπερέντασης και απόστασης πρέπει να
      διακρίνουν με ασφάλεια το ρεύμα σφάλματος από το ρεύμα φορτίου και να απομονώνουν
      επιλεκτικά μόνο το στοιχείο με το σφάλμα,
\item \textbf{τη σταθερότητα της τάσης}, γιατί σε ασθενές δίκτυο οι χειρισμοί πυκνωτών και οι
      μεταβολές του φορτίου προκαλούν μεγάλα βήματα τάσης και οι βυθίσεις από τα σφάλματα
      εκτείνονται σε μεγαλύτερη περιοχή,
\item \textbf{την ποιότητα ισχύος}, γιατί οι αρμονικές και οι διακυμάνσεις της τάσης που
      προκαλούν τα φορτία είναι τόσο μικρότερες όσο μεγαλύτερη είναι η ισχύς βραχυκύκλωσης, και
\item \textbf{την ευσταθή λειτουργία των μετατροπέων}, όπως εξηγείται παρακάτω.
\end{itemize}

Αντίθετα, η ισχύς βραχυκύκλωσης δεν πρέπει να υπερβαίνει την ικανότητα διακοπής των διακοπτών
ισχύος και την αντοχή του εξοπλισμού, οπότε ο ΔΣΜ τη διατηρεί εντός ορίων.

Για τη σύνδεση μονάδων με μετατροπείς χρησιμοποιείται ο λόγος βραχυκύκλωσης (Short Circuit
Ratio, SCR)

\begin{equation}
\mathrm{SCR} = \frac{S_{\mathrm{sc}}}{P_n}
\label{eq:scr}
\end{equation}

όπου $P_n$ η ονομαστική ισχύς της μονάδας. Δίκτυο με SCR μικρότερο από 3 περίπου
χαρακτηρίζεται ασθενές και με SCR μικρότερο από 2 πολύ ασθενές. Σε ασθενές δίκτυο η τάση του
σημείου σύνδεσης μεταβάλλεται έντονα με το ρεύμα του ίδιου του μετατροπέα. Ο βρόχος κλειδωμένης
φάσης και οι ελεγκτές ενός μετατροπέα που ακολουθεί το δίκτυο αλληλεπιδρούν τότε με αυτό και
μπορεί να προκαλέσουν ταλαντώσεις ή απώλεια του συγχρονισμού του μετατροπέα. Τα φαινόμενα αυτά
κατατάσσονται στην ευστάθεια που οφείλεται στους μετατροπείς (converter-driven stability).

\emph{Παράδειγμα.} Ένας ζυγός 150 kV έχει ισχύ βραχυκύκλωσης 1\,500 MVA, δηλαδή, από την
\eqref{eq:ssc}, ρεύμα τριφασικού βραχυκυκλώματος 5,8 kA περίπου. Σε αυτόν συνδέεται αιολικό
πάρκο 300 MW, οπότε SCR = 5. Αν σε ώρες χαμηλής ζήτησης τεθούν εκτός λειτουργίας
οι σύγχρονες μονάδες της περιοχής και η ισχύς βραχυκύκλωσης μειωθεί στα 750 MVA, ο λόγος πέφτει στο 2,5 και το
δίκτυο γίνεται ασθενές για το πάρκο. Σύμφωνα με την \eqref{eq:dv-ssc}, η ίδια μεταβολή της
άεργης ισχύος του πάρκου μεταβάλλει τότε την τάση κατά το διπλάσιο.

Η μείωση της ισχύος βραχυκύκλωσης αντιμετωπίζεται με:

\begin{itemize}
\item σύγχρονους αντισταθμιστές, που αυξάνουν την ισχύ βραχυκύκλωσης και την αδράνεια,
\item διατήρηση σε λειτουργία σύγχρονων μονάδων σε κρίσιμες περιοχές, με κόστος για την
      κατανομή της παραγωγής,
\item μετατροπείς σχηματισμού δικτύου (grid-forming), οι οποίοι συμπεριφέρονται ως πηγές τάσης
      πίσω από σύνθετη αντίσταση, όπως οι σύγχρονες μηχανές, και αποκρίνονται ακαριαία στις
      μεταβολές της τάσης και της γωνίας του δικτύου, εντός όμως του ορίου ρεύματός τους, και
\item προσαρμογή των ρυθμίσεων της προστασίας στα μικρότερα ρεύματα σφάλματος.
\end{itemize}

\section{Επανεκκίνηση και λειτουργία σε νησίδα}
\label{sec:restoration}

Μετά από γενική διακοπή (blackout) σε ολόκληρο το σύστημα ή σε μεγάλο τμήμα του, ο ΔΣΜ
αποκαθιστά την τροφοδότηση σύμφωνα με το σχέδιο αποκατάστασης που προβλέπει ο Κανονισμός (ΕΕ)
2017/2196 για την κατάσταση έκτακτης ανάγκης και την αποκατάσταση (Network Code on Emergency and
Restoration, NC ER). Η επανατροφοδότηση ξεκινά είτε από γειτονικά συστήματα μέσω των
διασυνδέσεων, με στρατηγική από πάνω προς τα κάτω (top-down), είτε από μονάδες του ίδιου του
συστήματος με ικανότητα επανεκκίνησης, με στρατηγική από κάτω προς τα πάνω (bottom-up).

Ικανότητα επανεκκίνησης (black start capability) έχει μια μονάδα που μπορεί να εκκινήσει χωρίς
εξωτερική τροφοδότηση και να ηλεκτρίσει τμήμα του δικτύου. Παραδοσιακά την παρέχουν
υδροηλεκτρικοί σταθμοί, που χρειάζονται ελάχιστη ισχύ για τα βοηθητικά τους, και αεριοστρόβιλοι
με ντιζελογεννήτρια εκκίνησης· εξετάζονται επίσης σταθμοί αποθήκευσης με μετατροπείς
σχηματισμού δικτύου. Οι μονάδες επανεκκίνησης δοκιμάζονται περιοδικά. Στη στρατηγική από κάτω
προς τα πάνω:

\begin{enumerate}
\item οι μονάδες επανεκκίνησης ηλεκτρίζουν επιλεγμένες διαδρομές του δικτύου μεταφοράς προς
      μεγάλους σταθμούς, ώστε να τροφοδοτηθούν τα βοηθητικά τους και να επανεκκινήσουν,
\item σχηματίζονται νησίδες, στις οποίες η παραγωγή και το φορτίο αυξάνονται σταδιακά, και
\item οι νησίδες συγχρονίζονται μεταξύ τους και με τα γειτονικά συστήματα.
\end{enumerate}

Τα τεχνικά προβλήματα της αποκατάστασης συνδέονται άμεσα με τα θέματα του κεφαλαίου:

\begin{itemize}
\item \textbf{Υπερτάσεις.} Οι γραμμές ηλεκτρίζονται χωρίς φορτίο και παράγουν άεργη ισχύ, που
      αυξάνει την τάση λόγω του φαινομένου Ferranti (Ενότητα «\nameref{sec:sil}»). Οι λίγες
      μονάδες σε λειτουργία πρέπει να την απορροφήσουν εντός του ορίου υποδιέγερσης, γι' αυτό
      συνδέονται πηνία, ηλεκτρίζονται σύντομες διαδρομές και συνδέεται νωρίς φορτίο. Η ζεύξη
      μετασχηματιστών χωρίς φορτίο προκαλεί επιπλέον ρεύματα εκκίνησης (inrush) και προσωρινές
      υπερτάσεις λόγω αρμονικού συντονισμού.
\item \textbf{Ρύθμιση συχνότητας.} Η νησίδα έχει μικρή αδράνεια και λίγες μονάδες, οπότε το
      φορτίο συνδέεται σε μικρά βήματα, ώστε η μεταβολή της συχνότητας να παραμένει
      περιορισμένη. Τα φορτία που επανατροφοδοτούνται μετά από μακρά διακοπή απορροφούν αρχικά
      περισσότερη ισχύ από την κανονική (cold load pickup), γιατί τα θερμοστατικά φορτία
      λειτουργούν τότε όλα ταυτόχρονα.
\item \textbf{Ασθενές δίκτυο.} Η ισχύς βραχυκύκλωσης είναι μικρή, οπότε η προστασία πρέπει να
      λειτουργεί με μικρά ρεύματα σφάλματος, και οι μονάδες με μετατροπείς που ακολουθούν το
      δίκτυο συνδέονται μόνο όταν η νησίδα αποκτήσει επαρκή ισχύ βραχυκύκλωσης.
\item \textbf{Συγχρονισμός.} Για τη σύνδεση δύο νησίδων ή το κλείσιμο ενός βρόχου, οι τάσεις
      στις δύο πλευρές του διακόπτη πρέπει να έχουν παραπλήσιο μέτρο, συχνότητα και γωνία.
      Μεγάλη διαφορά γωνίας σε ανοιχτό βρόχο (standing phase angle) προκαλεί κατά το κλείσιμο
      μεγάλη ροή ισχύος και καταπόνηση των μονάδων και περιορίζεται με ανακατανομή της παραγωγής
      πριν από τον χειρισμό.
\end{itemize}

Η ικανότητα λειτουργίας σε νησίδα (island operation capability) απαιτεί από μια μονάδα να
ρυθμίζει τη συχνότητα και την τάση της νησίδας και να παραμένει σε ευσταθή λειτουργία μετά από
τις μεταβολές του φορτίου που συνοδεύουν τον σχηματισμό της. Ο ρυθμιστής στροφών πρέπει να
ανιχνεύει τη μετάβαση σε λειτουργία νησίδας και να αλλάζει τρόπο ρύθμισης, για παράδειγμα σε
ισόχρονη ρύθμιση χωρίς στατισμό. Συγγενής ικανότητα είναι η ταχεία επανασύνδεση (quick
re-synchronisation): μια μονάδα που αποσυνδέεται από το δίκτυο παραμένει σε λειτουργία
τροφοδοτώντας μόνο τα βοηθητικά της (houseload operation), ώστε να επανασυνδεθεί αμέσως μόλις
επανέλθει η τάση. Ο RfG προβλέπει τις τρεις ικανότητες για μονάδες τύπου C και D· η ικανότητα
επανεκκίνησης δεν είναι υποχρεωτική και παρέχεται κατόπιν αιτήματος του ΔΣΜ.

\section{Ερωτήσεις και ασκήσεις}

\begin{enumerate}
\item \textbf{SVC και STATCOM.} Ένας SVC και ένας STATCOM έχουν ονομαστική χωρητική ισχύ
      100 Mvar. Πόση άεργη ισχύ αποδίδει ο καθένας όταν η τάση του ζυγού πέσει στα 0,7 p.u.;
      \emph{Απάντηση:} 49 Mvar και 70 Mvar αντίστοιχα.
\end{enumerate}

% TODO: παραδείγματα και λυμένες ασκήσεις (σελίδα 02b-paradeigmata.md, όπως στο Κεφάλαιο 1)

% ---------------------------------------------------------------------------
% BACKUP: εκτός δομής, για μελλοντική χρήση.
% Μελέτη περίπτωσης: γενική διακοπή στην Ιβηρική Χερσόνησο, 28 Απριλίου 2025.
% Η τελική έκθεση της ομάδας εμπειρογνωμόνων του ENTSO-E (Μάρτιος 2026) αποδίδει
% το συμβάν κυρίως σε υπερτάσεις και ανεπαρκή ρύθμιση τάσης, σε συνδυασμό με
% ταλαντώσεις που προηγήθηκαν. Πιθανή θέση: πλαίσιο μετά την ενότητα ρύθμισης
% τάσης ή στην ενότητα επανεκκίνησης.
% ---------------------------------------------------------------------------
